% -----------------------------------------------------------------------------
% Project.tex
%
% Example for the CalPolySeniorProject LaTeX class.
% - Includes: 2 figures, 1 table, and 1 equation (minimum).
% - No external figure files required: figures use boxed placeholders.
%
% Replace the metadata setters and content with your own project material.
% -----------------------------------------------------------------------------

\documentclass[
    % Optional margin overrides (defaults match the senior project format).
    left=1.5in,
    right=1.25in,
    top=1.25in,
    bottom=1.25in
]{CalPolySeniorProject}

% -----------------------------------------------------------------------------
% Metadata (fill these in for your report)
% -----------------------------------------------------------------------------
\setprojecttitle{Telemetry for Safety-Critical Railway Systems}
\setauthor{1}{Claude E. Shannon}{claude.shannon@calpoly.edu}
\setauthor{2}{Edith Clarke}{edith.clarke@calpoly.edu}
\setprofessor{Dr.\ Gordon E.\ Moore}
\setclass{EE 600}
\setdepartment{Electrical Engineering Department}
\setquarter{Fall 2025}

% Bibliography file for this example.
\addbibresource{Sources.bib}

\begin{document}

% -----------------------------------------------------------------------------
% Title page (provided by the class)
% NOTE: The class expects Images/Logo.pdf to exist.
% -----------------------------------------------------------------------------
\maketitlepage

% -----------------------------------------------------------------------------
% Front matter
% -----------------------------------------------------------------------------
\pagenumbering{roman}
\pagestyle{plain}

\tableofcontents

\clearpage
\listoffigures
\addcontentsline{toc}{section}{List of Figures}

\clearpage
\listoftables
\addcontentsline{toc}{section}{List of Tables}

% -----------------------------------------------------------------------------
% Abstract
% -----------------------------------------------------------------------------
\section*{Abstract}
\addcontentsline{toc}{section}{Abstract}

This document demonstrates the CalPolySeniorProject class using anonymized
metadata and placeholder content themed around electrical engineering. The
example includes figures, a table, and an equation to verify common report
features. Replace the text, citations, and placeholders with your real system
architecture, analysis, and results.

% -----------------------------------------------------------------------------
% Main matter
% -----------------------------------------------------------------------------
\clearpage
\pagenumbering{arabic}
\setcounter{page}{1}

\section{Introduction}

Electrical engineering designs succeed or fail on constraints: bandwidth,
power, timing, reliability, and implementation complexity. A good report makes
those constraints explicit and then shows how each design choice trades one
resource for another.

This example uses \texttt{biblatex} citations. For instance, Shannon's classic
result connects capacity, bandwidth, and signal-to-noise ratio (SNR)
\autocite{shannon1948}.

\section{Background}

In practical wireless links, packet delivery is rarely a step function of
distance. Instead, there is often a ``transitional region'' where delivery
smoothly degrades as the receiver moves away from the transmitter
\autocite{zuniga2004transitional}.

\subsection{A Small, On-Theme Equation}

A common sanity check is whether the required throughput is even plausible.
The (idealized) Shannon capacity is:

\begin{equation}
    C = B \log_2\!\left(1 + \mathrm{SNR}\right),
    \label{eq:shannon_capacity}
\end{equation}

where \(C\) is channel capacity (bit/s) and \(B\) is bandwidth (Hz).

\section{System Overview}

Figure~\ref{fig:arch} is a placeholder for an architecture diagram. Replace
the boxed region with a real figure (or keep the box if your reviewers love
minimalism).

\begin{figure}[!ht]
    \centering
    \fbox{%
        \parbox[c][1.75in][c]{0.90\linewidth}{%
            \centering
            \textbf{Placeholder Figure 1: System Architecture Diagram}\par
            \vspace{0.5em}
            Replace this box with your real figure (PDF/PNG/JPG).%
        }%
    }
    \caption[Architecture diagram placeholder]{Example architecture diagram
    placeholder.}
    \label{fig:arch}
\end{figure}

\section{Results}

Table~\ref{tab:results} demonstrates a simple results table. Use this pattern
for simulation sweeps, lab measurements, or performance comparisons.

\begin{table}[!ht]
    \centering
    \caption[Example results table]{Example performance summary for three
    configurations.}
    \vspace{0.5\baselineskip}
    \begin{tabular}{lccc}
        \toprule
        \textbf{Config} &
        \textbf{Reliability} &
        \textbf{Latency [ms]} &
        \textbf{Energy [mJ]} \\
        \midrule
        Baseline & 0.82 & 140 & 9.5 \\
        Tuned    & 0.93 & 110 & 8.1 \\
        Robust   & 0.97 & 125 & 9.0 \\
        \bottomrule
    \end{tabular}
    \label{tab:results}
\end{table}

Figure~\ref{fig:trend} is a placeholder for a results plot. Replace with your
actual plot export (recommended: PDF for vector quality).

\begin{figure}[!ht]
    \centering
    \fbox{%
        \parbox[c][1.75in][c]{0.90\linewidth}{%
            \centering
            \textbf{Placeholder Figure 2: Results Plot}\par
            \vspace{0.5em}
            Replace this box with your reliability/latency plot.%
        }%
    }
    \caption[Results plot placeholder]{Example results plot placeholder.}
    \label{fig:trend}
\end{figure}

\section{Conclusion}

This example file exists to prove the template is usable for a real report: it
compiles cleanly, produces a title page, prints a TOC/LOF/LOT, supports
citations, and includes typical technical elements. Replace the placeholders
with your system design, assumptions, validation, and final conclusions.

% -----------------------------------------------------------------------------
% Bibliography
% -----------------------------------------------------------------------------
\clearpage
\section{Works Cited}
\printbibliography[heading=none]

% -----------------------------------------------------------------------------
% Appendices
% -----------------------------------------------------------------------------
\beginappendices

\appsubsection{Example Appendix}

This appendix entry demonstrates the class appendix helper. Add more appendices
by calling \verb|\appsubsection{Example}| again.

\end{document}
